% USENIX-style report draft for INFO5995 Assignment 1
\documentclass[letterpaper,twocolumn,10pt]{article}
\usepackage{usenix,epsfig,endnotes}
\begin{document}

\date{}

\title{\Large \bf Predictable Session Token Generation in an Android APK Due to \texttt{java.util.Random}}

\author{
{\rm Rongbang Cheng}\\
INFO5995 Assignment 1 Group
\and
{\rm Team Members}\\
INFO5995
}

\maketitle
\thispagestyle{empty}

\subsection*{Abstract}
We analyzed the provided Android APK and selected one core vulnerability that fits the Assignment 1 randomness/cryptography scope. After successful credential verification, the app generates a value named \texttt{sessionToken} using \texttt{java.util.Random}, stores it in \texttt{SharedPreferences}, and uses it as session-related state. Because session tokens require unpredictability, using a deterministic non-cryptographic PRNG is unsuitable for this security-sensitive role. Our evidence supports a bounded predictability-risk claim under realistic local-analysis conditions rather than a guaranteed authentication bypass claim.

\section{Introduction}

The provided APK implements a simple local registration and login flow using three main activities: \texttt{MainActivity}, \texttt{Login}, and \texttt{Profile}. In the observed flow, \texttt{MainActivity} stores user credentials locally, \texttt{Login} checks the supplied credentials, and a successful login creates session-related state before opening \texttt{Profile}. We selected one core issue only: weak generation of the authentication-related \texttt{sessionToken} in \texttt{Login.generateSessionToken()}.

This issue is directly within the Assignment 1 scope because it is a misuse of randomness in a security-sensitive context, not merely a cosmetic random value. Our goal is therefore not to claim every possible weakness in the APK, but to present one defensible randomness vulnerability with clear code evidence, a bounded attacker model, and a realistic mitigation.

\section{System and Threat Model}

The relevant system components are \texttt{Login}, which verifies credentials and creates session state, and \texttt{Profile}, which represents the post-login state and clears the token on logout. The primary protected asset is the unpredictability and integrity of authentication-related session state represented by \texttt{sessionToken}. This value crosses an important trust boundary when it is generated and then persisted in \texttt{SharedPreferences("SessionPrefs")}.

Our attacker model is intentionally bounded. We assume a local-analysis attacker using an emulator, rooted device, or instrumented test environment who can estimate when a victim logged in and may be able to inspect or test locally stored token state. We do not assume backend compromise, remote code execution, or a proven downstream token-validation sink in the current build. This model is realistic for the assignment setting and avoids exaggerating impact beyond the available evidence.

\section{Vulnerability Discovery and Explanation}

Static analysis of the decompiled APK identified two visible uses of randomness. The first is \texttt{MainActivity.randomNumberGenerator()}, which produces a UI-style random integer and has no clear authentication role. The second is \texttt{Login.generateSessionToken()}, which creates a 16-character token using \texttt{new Random()} and repeated \texttt{nextInt(62)} calls. We selected the second path because the generated value is explicitly named \texttt{sessionToken}, created immediately after successful login, stored as session state, and cleared on logout.

This makes the token a security-sensitive value rather than a display-only value. \texttt{java.util.Random} is a deterministic PRNG rather than a cryptographically secure random generator. Session tokens require unpredictability, so using \texttt{Random} for this purpose is an inappropriate design choice. The security problem is therefore not ``the app uses randomness,'' but that it uses a predictable generator to create authentication-related state.

\section{Attack Scenario and Impact}

In the observed code path, successful credential verification triggers \texttt{createSession()}, which stores a newly generated \texttt{sessionToken}, and then transitions the user to \texttt{Profile}. Under a realistic local-analysis setup, an attacker who can narrow the login timing window can reproduce the same token-generation logic offline and derive candidate outputs from nearby PRNG seeds. If the attacker can also inspect or test token state in that environment, they can compare candidate outputs against the observed token or token-dependent behavior.

The strongest defensible impact is weakened unpredictability of authentication-related session state. Our current evidence shows token generation, storage, retrieval, and clearing more clearly than downstream token validation. Therefore, we do not claim guaranteed authentication bypass, session hijack, or remote compromise. Instead, we claim that the app creates session-related state with a predictable PRNG where a cryptographically secure generator should have been used.

\section{Fix and Mitigation}

The primary mitigation is to replace \texttt{java.util.Random} with \texttt{SecureRandom} when generating \texttt{sessionToken}. This directly improves unpredictability for security-sensitive token creation. In addition, the app should make token validation explicit in protected flow and enforce clear invalidation or rotation semantics, so that authentication state is both unpredictable and correctly enforced throughout its lifecycle.

\section{Tooling and AI Use}

We used APK decompilation and static code inspection to locate activity structure, authentication flow, and randomness usage. AI assistance was used to help organize evidence, compare candidate findings against the rubric, and rehearse bounded tutorial Q\&A answers. All final claims were checked against decompiled code locations and screenshot evidence before inclusion in the report.

\end{document}
